% Source: https://artofproblemsolving.com/community/c3h2029806_happy_pi_day_maybe_a_competition
% Stable identity: aops-topic-2029806-q13
Let the roots be integers \(r,s\) (possibly equal). By Vieta's formulas \(r+s=-m\) and \(rs=4m\), so \(rs=-4(r+s)\), i.e.
\[
(r+4)(s+4)=16 .
\]
Conversely, any integers \(r,s\) with this property give \(m=-(r+s)\) and the quadratic \((x-r)(x-s)\) has constant term \(rs=4m\).

The unordered factorisations \(16=uv\) are \(1\cdot16,\ 2\cdot8,\ 4\cdot4,\ (-1)(-16),\ (-2)(-8),\ (-4)(-4)\), with \(u+v=17,10,8,-17,-10,-8\). Since \(m=-(r+s)=8-(u+v)\), the possible values are
\[
m=-9,\ -2,\ 0,\ 25,\ 18,\ 16,
\]
all distinct. Their mean is \(\dfrac{-9-2+0+25+18+16}{6}=\dfrac{48}{6}=\boxed{8}\).
